\documentclass[11pt,a4paper]{article}

\usepackage[margin=2.2cm]{geometry}
\usepackage{kotex}
\usepackage{amsmath,amssymb}
\usepackage{booktabs}
\usepackage{array}
\usepackage{xcolor}
\usepackage{listings}
\usepackage[most]{tcolorbox}
\usepackage{enumitem}
\usepackage{hyperref}
\usepackage{fancyhdr}
\usepackage{titlesec}

\definecolor{codebg}{RGB}{247,247,247}
\definecolor{codeframe}{RGB}{210,210,210}
\definecolor{keywordcolor}{RGB}{0,70,160}
\definecolor{commentcolor}{RGB}{0,120,80}
\definecolor{stringcolor}{RGB}{170,50,50}
\definecolor{titleblue}{RGB}{35,75,130}

\hypersetup{colorlinks=true,linkcolor=titleblue,urlcolor=titleblue}

\pagestyle{fancy}
\fancyhf{}
\lhead{C++ Programming}
\rhead{Day 13}
\cfoot{\thepage}

\titleformat{\section}{\Large\bfseries\color{titleblue}}{\thesection.}{0.6em}{}
\titleformat{\subsection}{\large\bfseries}{\thesubsection.}{0.5em}{}

\lstdefinestyle{cppstyle}{
    language=C++,
    backgroundcolor=\color{codebg},
    frame=single,
    rulecolor=\color{codeframe},
    basicstyle=\ttfamily\small,
    keywordstyle=\color{keywordcolor}\bfseries,
    commentstyle=\color{commentcolor},
    stringstyle=\color{stringcolor},
    numbers=left,
    numberstyle=\tiny\color{gray},
    numbersep=8pt,
    tabsize=4,
    showstringspaces=false,
    breaklines=true,
    columns=fullflexible
}
\lstset{style=cppstyle}

\newtcolorbox{learningbox}{colback=blue!4,colframe=titleblue,title=\textbf{학습 목표},fonttitle=\bfseries}
\newtcolorbox{importantbox}{colback=yellow!8,colframe=orange!70!black,title=\textbf{핵심 포인트},fonttitle=\bfseries}
\newtcolorbox{practicebox}{colback=red!3,colframe=red!55!black,title=\textbf{실습},fonttitle=\bfseries}

\begin{document}

\begin{titlepage}
    \centering
    \vspace*{3cm}
    {\Huge\bfseries C++ Programming\par}
    \vspace{0.8cm}
    {\LARGE Day 13: File I/O와 Exception Handling\par}
    \vspace{2cm}
    {\large ofstream, ifstream, ios::app, File Parsing, try/catch, throw\par}
    \vfill
\end{titlepage}

\tableofcontents
\newpage

\section{학습 목표}
\begin{learningbox}
이 장을 마치면 다음 내용을 설명하고 직접 코드로 작성할 수 있다.
\begin{itemize}[leftmargin=1.5em]
    \item \texttt{ofstream}을 이용해 파일에 데이터를 저장하기
    \item \texttt{ifstream}을 이용해 파일의 데이터를 읽기
    \item \texttt{ios::app}을 이용해 기존 파일 끝에 데이터를 추가하기
    \item \texttt{getline()}과 스트림 추출 연산자 \texttt{>>}의 차이 이해하기
    \item 파일의 데이터를 여러 변수로 나누어 읽는 File Parsing 수행하기
    \item \texttt{try}와 \texttt{catch}를 이용해 예외 처리 구조 작성하기
    \item \texttt{throw}와 \texttt{runtime\_error}를 이용해 예외 발생시키기
    \item \texttt{exception}과 \texttt{what()}을 이용해 예외 메시지 처리하기
\end{itemize}
\end{learningbox}

\section{Day 13 개요}
Day 13에서는 프로그램 외부의 파일과 데이터를 주고받는 File I/O와 실행 중 발생할 수 있는 문제를 처리하는 Exception Handling을 학습한다.
파일 출력에서는 \texttt{ofstream}, 파일 입력에서는 \texttt{ifstream}을 사용하고,
기존 파일의 내용을 유지하면서 데이터를 추가할 때는 \texttt{ios::app}을 사용한다.
또한 파일의 데이터를 변수로 분리하여 처리하는 File Parsing을 학습한다.
마지막으로 \texttt{try}, \texttt{catch}, \texttt{throw}를 이용한 기본적인 예외 처리 방법을 다룬다.

\begin{center}
\begin{tabular}{ll}
\toprule
주제 & 핵심 역할 \\
\midrule
\texttt{ofstream} & 파일에 데이터 쓰기 \\
\texttt{ifstream} & 파일에서 데이터 읽기 \\
\texttt{ios::app} & 기존 파일 끝에 데이터 추가 \\
File Parsing & 파일 데이터를 변수로 분리 \\
\texttt{try / catch} & 예외 발생 가능 코드와 처리 코드 분리 \\
\texttt{throw} & 예외 발생시키기 \\
\bottomrule
\end{tabular}
\end{center}

\begin{importantbox}
Day 13의 핵심은 파일 스트림을 이용해 외부 데이터를 읽고 쓰는 흐름과,
예외가 발생했을 때 프로그램의 제어 흐름이 \texttt{catch}로 이동하는 구조를 이해하는 것이다.
\end{importantbox}

\section{File Output: ofstream}
\texttt{ofstream}은 output file stream의 의미로, 프로그램에서 파일에 데이터를 출력할 때 사용한다.
파일 입출력 기능을 사용하려면 \texttt{<fstream>} 헤더가 필요하다.

\subsection{파일 열기와 데이터 쓰기}
다음 코드는 \texttt{data.txt} 파일을 연다.

\begin{lstlisting}
ofstream file("data.txt");
\end{lstlisting}

파일이 존재하지 않으면 새 파일을 만들 수 있다.
파일에 데이터를 쓸 때는 \texttt{cout}과 마찬가지로 \texttt{<<} 연산자를 사용한다.

\begin{lstlisting}
file << "Hello" << endl;
file << 100 << endl;
\end{lstlisting}

파일이 정상적으로 열렸는지는 \texttt{is\_open()}으로 확인할 수 있다.

\begin{lstlisting}
if (!file.is_open())
{
    cout << "Failed to open file." << endl;
    return 1;
}
\end{lstlisting}

\texttt{main()}에서 \texttt{return 0;}은 정상 종료를 의미하고,
0이 아닌 값은 일반적으로 오류가 발생한 종료 상태를 나타낸다.
파일 사용이 끝나면 \texttt{close()}로 닫는다.

\subsection{최종 코드}
\begin{lstlisting}
#include <iostream>
#include <fstream>
#include <string>

using namespace std;

int main()
{
    ofstream file("data.txt");

    if (!file.is_open())
    {
        cout << "Failed to open file." << endl;
        return 1;
    }

    string name = "Alice";
    int score = 95;

    file << "Name: " << name << endl;
    file << "Score: " << score << endl;

    file.close();

    cout << "File saved successfully." << endl;

    return 0;
}
\end{lstlisting}

파일에는 다음 내용이 저장된다.

\begin{verbatim}
Name: Alice
Score: 95
\end{verbatim}

\subsection{실습}
\begin{practicebox}
\texttt{student.txt} 파일을 생성하고 Bob의 이름, 전공, 점수를 저장한다.
파일이 정상적으로 열렸는지 확인하고 저장이 끝나면 파일을 닫는다.
\end{practicebox}

\begin{lstlisting}
#include <iostream>
#include <fstream>
#include <string>

using namespace std;

int main()
{
    ofstream file("student.txt");

    if (!file.is_open())
    {
        cout << "Failed to open file." << endl;
        return 1;
    }

    string name = "Bob";
    string major = "Physics";
    int score = 88;

    file << "Name: " << name << endl;
    file << "Major: " << major << endl;
    file << "Score: " << score << endl;

    file.close();

    cout << "Student data saved " << endl;

    return 0;
}
\end{lstlisting}

\section{File Input: ifstream}
\texttt{ifstream}은 input file stream의 의미로, 파일에 저장된 데이터를 프로그램으로 읽어올 때 사용한다.
\texttt{cin}이 키보드 입력을 받는다면 \texttt{ifstream}은 파일에서 입력을 받는다.

\subsection{파일 열기}
다음과 같이 읽을 파일을 지정한다.

\begin{lstlisting}
ifstream file("data.txt");
\end{lstlisting}

\texttt{ifstream}은 읽을 파일이 존재해야 하므로 파일이 정상적으로 열렸는지 확인하는 것이 중요하다.

\subsection{getline으로 한 줄 읽기}
\texttt{>>} 연산자는 공백을 기준으로 값을 나누어 읽는다.
공백을 포함한 한 줄 전체를 읽으려면 \texttt{getline()}을 사용한다.

\begin{lstlisting}
string line;
getline(file, line);
\end{lstlisting}

파일의 모든 줄을 읽을 때는 다음 패턴을 사용할 수 있다.

\begin{lstlisting}
while (getline(file, line))
{
    cout << line << endl;
}
\end{lstlisting}

\texttt{getline()}이 한 줄을 읽는 데 성공하는 동안 반복하고,
더 이상 읽을 데이터가 없으면 반복문을 종료한다.

\subsection{최종 코드}
\begin{lstlisting}
#include <iostream>
#include <fstream>
#include <string>

using namespace std;

int main()
{
    ifstream file("data.txt");

    if (!file.is_open())
    {
        cout << "Failed to open file." << endl;
        return 1;
    }

    string line;

    while (getline(file, line))
    {
        cout << line << endl;
    }

    file.close();

    return 0;
}
\end{lstlisting}

\subsection{실습}
\begin{practicebox}
\texttt{student.txt}를 \texttt{ifstream}으로 열고,
\texttt{getline()}과 \texttt{while}을 사용하여 파일의 모든 줄을 출력한다.
\end{practicebox}

\begin{lstlisting}
#include <iostream>
#include <fstream>
#include <string>

using namespace std;

int main()
{
    ifstream file("student.txt");

    if (!file.is_open())
    {
        cout << "Failed to open file." << endl;
        return 1;
    }

    string line;

    while (getline(file, line))
    {
        cout << line << endl;
    }

    file.close();

    return 0;
}
\end{lstlisting}

\section{Append Mode: ios::app}
기본 \texttt{ofstream}으로 기존 파일을 열면 이전 내용을 덮어쓰게 된다.
기존 내용을 유지하면서 파일 끝에 새로운 데이터를 추가하려면 \texttt{ios::app}을 사용한다.

\subsection{기본 구조}
\begin{lstlisting}
ofstream file("student.txt", ios::app);
\end{lstlisting}

\texttt{app}은 append를 의미하며, 이후 출력되는 내용은 기존 파일의 끝에 추가된다.

\begin{center}
\begin{tabular}{ll}
\toprule
형태 & 동작 \\
\midrule
\texttt{ofstream file("data.txt")} & 기존 내용을 덮어쓰며 출력 \\
\texttt{ofstream file("data.txt", ios::app)} & 기존 내용 뒤에 추가 \\
\bottomrule
\end{tabular}
\end{center}

\subsection{최종 코드}
\begin{lstlisting}
#include <iostream>
#include <fstream>
#include <string>

using namespace std;

int main()
{
    ofstream file("student.txt", ios::app);

    if (!file.is_open())
    {
        cout << "Failed to open file." << endl;
        return 1;
    }

    string name = "Alice";
    string major = "Mathematics";
    int score = 95;

    file << endl;
    file << "Name: " << name << endl;
    file << "Major: " << major << endl;
    file << "Score: " << score << endl;

    file.close();

    cout << "Student data appended." << endl;

    return 0;
}
\end{lstlisting}

\subsection{실습}
\begin{practicebox}
기존 \texttt{student.txt}에 Charlie의 정보를 추가한다.
\texttt{ios::app}을 사용하고 기존 학생과 구분되도록 빈 줄 하나를 먼저 출력한다.
\end{practicebox}

\begin{lstlisting}
#include <iostream>
#include <fstream>
#include <string>

using namespace std;

int main()
{
    ofstream file("student.txt", ios::app);

    if (!file.is_open())
    {
        cout << "Failed to open file." << endl;
        return 1;
    }

    file << endl;

    string name = "Charlie";
    string major = "ComputerScience";
    int score = 91;

    file << "Name: " << name << endl;
    file << "Major: " << major << endl;
    file << "Score: " << score << endl;

    file.close();

    cout << "Student data appended." << endl;

    return 0;
}
\end{lstlisting}

\section{File Parsing}
File Parsing은 파일에 저장된 데이터를 단순한 한 줄 문자열로 읽는 것이 아니라,
각 데이터를 변수에 나누어 저장하여 프로그램에서 사용할 수 있도록 처리하는 것이다.

\subsection{스트림 추출 연산자로 데이터 분리}
파일의 내용이 다음과 같다고 하자.

\begin{verbatim}
Alice 95
Bob 88
Charlie 91
\end{verbatim}

다음 코드로 이름과 점수를 각각의 변수에 읽을 수 있다.

\begin{lstlisting}
string name;
int score;

file >> name >> score;
\end{lstlisting}

여러 줄을 계속 읽으려면 입력 작업 자체를 \texttt{while}의 조건으로 사용할 수 있다.

\begin{lstlisting}
while (file >> name >> score)
{
    cout << name << " " << score << endl;
}
\end{lstlisting}

\subsection{읽은 데이터 계산하기}
읽은 점수를 합계와 개수에 반영할 수 있다.

\begin{lstlisting}
int total = 0;
int count = 0;

while (file >> name >> score)
{
    total += score;
    count++;
}
\end{lstlisting}

평균을 계산할 때 정수 나눗셈을 피하려면 \texttt{static\_cast<double>}을 사용한다.

\begin{lstlisting}
double average = static_cast<double>(total) / count;
\end{lstlisting}

\subsection{최종 코드}
\begin{lstlisting}
#include <iostream>
#include <fstream>
#include <string>

using namespace std;

int main()
{
    ifstream file("scores.txt");

    if (!file.is_open())
    {
        cout << "Failed to open file." << endl;
        return 1;
    }

    string name;
    int score;

    int total = 0;
    int count = 0;

    while (file >> name >> score)
    {
        cout << "Name: " << name
             << ", Score: " << score << endl;

        total += score;
        count++;
    }

    file.close();

    if (count > 0)
    {
        double average = static_cast<double>(total) / count;

        cout << "Total: " << total << endl;
        cout << "Average: " << average << endl;
    }

    return 0;
}
\end{lstlisting}

\subsection{실습}
\begin{practicebox}
\texttt{grades.txt}에 저장된 이름과 점수를 읽어 각 학생의 점수를 출력하고,
전체 점수의 총합과 평균을 계산한다.
\end{practicebox}

\begin{verbatim}
Tom 80
Jane 92
Mike 75
Emma 98
\end{verbatim}

\begin{lstlisting}
#include <iostream>
#include <fstream>
#include <string>

using namespace std;

int main()
{
    ifstream file("grades.txt");

    if (!file.is_open())
    {
        cout << "Failed to open file" << endl;
        return 1;
    }

    string name;
    int score;

    int total = 0;
    int count = 0;

    while (file >> name >> score)
    {
        cout << name << ": " << score << endl;
        total += score;
        count++;
    }

    file.close();

    if (count > 0)
    {
        double average = static_cast<double>(total) / count;

        cout << "Total: " << total << endl;
        cout << "Average: " << average << endl;
    }

    return 0;
}
\end{lstlisting}

\section{Exception Handling: try와 catch}
예외 처리는 프로그램 실행 중 발생할 수 있는 문제를 처리하는 방법이다.
\texttt{try} 블록에는 예외가 발생할 수 있는 코드를 작성하고,
예외가 발생하면 대응되는 \texttt{catch} 블록에서 처리한다.

\subsection{기본 구조}
\begin{lstlisting}
try
{
    throw 1;
}
catch (...)
{
    cout << "Error occurred." << endl;
}
\end{lstlisting}

\texttt{catch (...)}는 예외의 종류와 관계없이 예외를 잡는다.
\texttt{throw}가 실행되면 현재 \texttt{try} 블록의 나머지 코드는 실행되지 않고,
처리할 수 있는 \texttt{catch} 블록으로 제어가 이동한다.

\subsection{최종 코드}
\begin{lstlisting}
#include <iostream>

using namespace std;

int main()
{
    try
    {
        cout << "Try block started." << endl;

        throw 1;

        cout << "Try block finished." << endl;
    }
    catch (...)
    {
        cout << "Exception caught." << endl;
    }

    cout << "Program continues." << endl;

    return 0;
}
\end{lstlisting}

출력 결과는 다음과 같다.

\begin{verbatim}
Try block started.
Exception caught.
Program continues.
\end{verbatim}

\subsection{실습}
\begin{practicebox}
\texttt{try} 블록에서 \texttt{Start}를 출력한 뒤 \texttt{throw 100;}으로 예외를 발생시킨다.
\texttt{catch (...)}에서 예외를 처리하고, 이후 프로그램이 계속 실행되는지 확인한다.
\end{practicebox}

\begin{lstlisting}
#include <iostream>

using namespace std;

int main()
{
    try
    {
        cout << "Start" << endl;
        throw 100;
        cout << "End" << endl;
    }
    catch (...)
    {
        cout << "Exception detected." << endl;
    }

    cout << "Program finished." << endl;

    return 0;
}
\end{lstlisting}

\section{throw와 Standard Exception}
단순한 정수도 \texttt{throw}할 수 있지만,
실제 오류 정보를 전달하려면 표준 예외 클래스를 사용할 수 있다.
Day 13에서는 \texttt{runtime\_error}와 \texttt{exception}을 사용한다.

\subsection{runtime\_error}
\texttt{runtime\_error}를 사용하려면 \texttt{<stdexcept>} 헤더가 필요하다.

\begin{lstlisting}
throw runtime_error("Invalid value.");
\end{lstlisting}

예외는 다음과 같이 받을 수 있다.

\begin{lstlisting}
catch (const exception& e)
{
    cout << e.what() << endl;
}
\end{lstlisting}

\texttt{const exception\&}는 예외 객체를 복사하지 않고 참조로 받으며,
\texttt{what()}은 예외 객체에 저장된 오류 메시지를 반환한다.
\texttt{runtime\_error}는 표준 \texttt{exception} 계열의 예외이므로
\texttt{catch (const exception\& e)}로 처리할 수 있다.

\subsection{함수에서 예외 발생시키기}
예외를 발견하는 코드와 처리하는 코드는 서로 다른 함수에 있을 수 있다.

\begin{lstlisting}
double divide(double a, double b)
{
    if (b == 0)
    {
        throw runtime_error("Cannot divide by zero.");
    }

    return a / b;
}
\end{lstlisting}

함수에서 발생한 예외는 호출한 쪽의 \texttt{catch}에서 처리할 수 있다.

\subsection{최종 코드}
\begin{lstlisting}
#include <iostream>
#include <stdexcept>

using namespace std;

double divide(double a, double b)
{
    if (b == 0)
    {
        throw runtime_error("Cannot divide by zero.");
    }

    return a / b;
}

int main()
{
    double a;
    double b;

    cout << "Enter two numbers: ";
    cin >> a >> b;

    try
    {
        double result = divide(a, b);

        cout << "Result: " << result << endl;
    }
    catch (const exception& e)
    {
        cout << "Error: " << e.what() << endl;
    }

    return 0;
}
\end{lstlisting}

\subsection{실습}
\begin{practicebox}
\texttt{checkScore(int score)} 함수를 만들고 점수가 0보다 작거나 100보다 크면
\texttt{runtime\_error}를 발생시킨다.
정상 범위의 점수라면 점수를 출력하고, 예외는 \texttt{main()}의 \texttt{catch}에서 처리한다.
\end{practicebox}

\begin{lstlisting}
#include <iostream>
#include <stdexcept>

using namespace std;

void checkScore(int score)
{
    if (score < 0 || score > 100)
    {
        throw runtime_error("Score must be between 0 and 100.");
    }

    cout << "Valid score: " << score << endl;
}

int main()
{
    int score;
    cout << "Enter a score: ";
    cin >> score;

    try
    {
        checkScore(score);
    }
    catch (const exception& e)
    {
        cout << "Error: " << e.what() << endl;
    }

    return 0;
}
\end{lstlisting}

\section{개념 연결}
Day 13의 각 기능은 이전에 배운 C++ 문법과 연결된다.

\begin{center}
\begin{tabular}{ll}
\toprule
이전 학습 내용 & Day 13에서의 연결 \\
\midrule
스트림 입출력 & \texttt{cin}, \texttt{cout}에서 파일 스트림으로 확장 \\
참조 & \texttt{catch (const exception\& e)} \\
클래스와 상속 & \texttt{runtime\_error}와 \texttt{exception} \\
함수 & 함수 내부에서 \texttt{throw}하고 호출한 쪽에서 처리 \\
형 변환 & 평균 계산에서 \texttt{static\_cast<double>} 사용 \\
\bottomrule
\end{tabular}
\end{center}

\begin{importantbox}
\texttt{ofstream}과 \texttt{ifstream}은 기존의 스트림 입출력 문법을 파일로 확장한 것이다.
예외 처리에서는 문제가 발생한 위치에서 \texttt{throw}하고,
처리할 위치에서 \texttt{catch}하여 프로그램의 오류 처리 흐름을 분리할 수 있다.
\end{importantbox}

\section{종합 정리}
\begin{importantbox}
\begin{itemize}[leftmargin=1.5em]
    \item \texttt{ofstream}은 파일에 데이터를 출력할 때 사용한다.
    \item \texttt{ifstream}은 파일의 데이터를 읽을 때 사용한다.
    \item 파일이 정상적으로 열렸는지는 \texttt{is\_open()}으로 확인할 수 있다.
    \item \texttt{getline()}은 공백을 포함하여 한 줄 전체를 읽는다.
    \item \texttt{ios::app}은 기존 파일 내용을 유지하면서 파일 끝에 데이터를 추가한다.
    \item \texttt{file >> variable} 형태를 이용하면 파일 데이터를 변수로 나누어 읽을 수 있다.
    \item 평균 계산에서 정수 나눗셈을 피하려면 \texttt{static\_cast<double>}을 사용할 수 있다.
    \item \texttt{try}에는 예외가 발생할 수 있는 코드를 작성하고 \texttt{catch}에서 예외를 처리한다.
    \item \texttt{throw}가 실행되면 현재 \texttt{try}의 나머지 코드는 실행되지 않는다.
    \item \texttt{runtime\_error}는 오류 메시지를 포함한 예외를 발생시키는 데 사용할 수 있다.
    \item \texttt{catch (const exception\& e)}로 표준 예외를 받고 \texttt{e.what()}으로 메시지를 확인할 수 있다.
\end{itemize}
\end{importantbox}

\section{실습 파일 구성}
\begin{practicebox}
Day 13 파일은 다음과 같이 관리한다.
\begin{itemize}[leftmargin=1.5em]
    \item \texttt{01\_file\_output.cpp}
    \item \texttt{02\_file\_input.cpp}
    \item \texttt{03\_file\_append.cpp}
    \item \texttt{04\_file\_parsing.cpp}
    \item \texttt{05\_exception\_basic.cpp}
    \item \texttt{06\_throw\_catch.cpp}
    \item \texttt{data.txt}
    \item \texttt{scores.txt}
    \item \texttt{practice/p01\_file\_output.cpp}
    \item \texttt{practice/p02\_file\_input.cpp}
    \item \texttt{practice/p03\_file\_append.cpp}
    \item \texttt{practice/p04\_file\_parsing.cpp}
    \item \texttt{practice/p05\_exception\_basic.cpp}
    \item \texttt{practice/p06\_throw\_catch.cpp}
    \item \texttt{practice/student.txt}
    \item \texttt{practice/grades.txt}
\end{itemize}
\end{practicebox}

\end{document}
